\section{Conclusion} 
Ce chapitre a présenté la démarche méthodologique mise en œuvre depuis l’installation des plots jusqu’au suivi des cultures, en passant par le pilotage de l’irrigation, la collecte des données, l’entretien et la lutte contre les ravageurs et les maladies. Cette approche a permis de mettre en place une base de données issue directement du terrain, regroupant ainsi les paramètres climatiques, hydriques et physico-chimiques nécessaires au calcul de l'évapotranspiration et au pilotage de l’irrigation. Ces données réelles constituent le socle du développement de l’application, destiné à automatiser le  calcul de l'évapotranspiration, et le pilotage de l’irrigation, tout en restant adaptée aux conditions de notre milieu d'étude.
